\begin{table}[htbp]
\centering
\small
\caption{Problemas B e C, Passo 4 (google/gemini-2.5-flash-lite): few-shot com exemplos rotulados pela métrica do A (H2).}
\label{tab:bloco1_BC_passo4_exemplos_da_metrica_flashlite}
\begin{tabular}{lrrlrrrlr}
\toprule
 & \multicolumn{4}{c}{Problema B} & \multicolumn{4}{c}{Problema C} \\ \cmidrule(lr){2-5} \cmidrule(lr){6-9}
Regime & Consist. & $\kappa$ & Faixa & $\Delta$ & Consist. & $\kappa$ & Faixa & $\Delta$ \\
\midrule
zero-shot & 72,3 (68,0–76,0) & 0,258 & razoável & +0,0 & 92,7 (92,0–94,0) & 0,852 & quase perfeita & +0,0 \\
4-shot & 86,1 (76,0–95,8) & 0,659 & substancial & +13,8 & 93,1 (92,7–93,8) & 0,861 & quase perfeita & +0,4 \\
10-shot & 94,1 (91,1–97,8) & 0,881 & quase perfeita & +21,7 & 95,6 (93,3–97,8) & 0,911 & quase perfeita & +2,9 \\
20-shot & 94,6 (88,8–97,5) & 0,892 & quase perfeita & +22,3 & 92,9 (88,8–95,0) & 0,858 & quase perfeita & +0,2 \\
40-shot & 91,7 (86,7–98,3) & 0,832 & quase perfeita & +19,3 & 87,8 (78,3–95,0) & 0,751 & substancial & $-$4,9 \\
\bottomrule
\multicolumn{9}{p{0.95\linewidth}}{\footnotesize Problemas B e C, Passo 4 (hipótese H2): os k exemplos mostrados ao LLM em B e C são rotulados pela métrica zero-shot do A (seleção por margem, k/2 por classe) e o LLM classifica os demais pontos. Os exemplos few-shot de B e C são rotulados por $\hat{W}$ do A com os centróides do A; a consistência usa os centróides locais, recalculados pelos rótulos do LLM no próprio problema. Consist.\ = consistência (\%). $\kappa$ = Cohen's $\kappa$ (faixas de Landis e Koch, 1977). Médias das sementes (mín–máx); $\Delta$ em relação ao zero-shot. Perceptron, configuração base.} \\
\end{tabular}
\end{table}
